\section{Explicit Examples: Computing $A_\infty$ Operations and Cohomological Descent}
\label{sec:examples}

This section provides three fully worked examples demonstrating the $A_\infty$ chiral algebra framework. For each example, we:
\begin{enumerate}[label=(\roman*)]
\item Specify the field content and BV-BRST action;
\item Compute the differential $Q = m_1$ explicitly;
\item Derive the $\lambda$-bracket $m_2$ from the free propagator;
\item Determine higher operations $m_{k \geq 3}$ from Feynman diagrams (when nonzero);
\item Compute cohomology $H^\bullet(\A, Q)$ and verify PVA structure;
\item Provide explicit integral computations for at least one operation.
\end{enumerate}

These examples illustrate different phenomena:
\begin{itemize}
\item \textbf{Example \ref{ex:free_chiral}:} Free chiral multiplet---all $m_{k \geq 3} = 0$ (no interactions); trivial cohomology.
\item \textbf{Example \ref{ex:LG}:} Landau--Ginzburg cubic (first nontrivial $m_3$ from triangle diagram); truncation at $k=4$ by degree counting.
\item \textbf{Example \ref{ex:abelian_CS}:} Abelian CS gauge theory (boundary affine Kac--Moody algebra $\widehat{\mathfrak{u}(1)}_k$); connections to quantum groups.
\end{itemize}

\begin{remark}[Analytic hypotheses]
All computations in this section are conditional on Hypotheses~\ref{hyp:H1}--\ref{hyp:H3} and Assumption~\ref{ass:H1-H4}. In all three examples, (H1)--(H4) can be verified directly: the BV data are explicit and finite-dimensional, propagators are standard holomorphic--topological kernels with the required regularity, interaction vertices are polynomial, and factorization compatibility follows from the general results of \cite{GRW21,CDG20}. The free multiplet and abelian Chern--Simons examples are in fact exact at one loop.
\end{remark}

\subsection{Example I: Free Chiral Multiplet}
\label{ex:free_chiral}

The free chiral multiplet provides the simplest nontrivial example, where higher operations $m_{k \geq 3}$ vanish identically.

\subsubsection{Field Content and Action}

A free chiral multiplet on $\R \times \C$ consists of:
\begin{itemize}
\item A complex scalar $\phi : \R \times \C \to \C$ (the bottom component, in cohomological degree~$0$);
\item A fermionic field $\psi$ (the BV antifield, in cohomological degree~$1$).
\end{itemize}

In the BV formalism after gauge fixing, the action is
\[
S_{\text{free}} = \int_{\R \times \C} \langle \psi, (d_t + \bar{\partial})\phi \rangle \, dt \, dz \, d\bar{z},
\]
where $d_t + \bar{\partial}$ is the HT kinetic operator.

\subsubsection{Local Observables}

The algebra of local observables is the polynomial algebra
\[
\A_{\text{free}} = \C[\phi_n, \psi_n \mid n \geq 0],
\]
generated by derivatives
\[
\phi_n := \frac{1}{n!} \partial_z^n \phi, \quad \psi_n := \frac{1}{n!} \partial_z^n \psi,
\]
with cohomological degrees $|\phi_n| = 0$, $|\psi_n| = 1$.

\subsubsection{BV-BRST Differential: Explicit Formulas}

\begin{proposition}[Free BV-BRST Differential]
\label{prop:free_BRST_detailed}
The BV-BRST differential $Q = \{S_{\text{free}}, -\}_{\mathrm{BV}}$ acts on generators by
\begin{align}
Q \cdot \phi_n &= \psi_n, \label{eq:Qphi_free}\\
Q \cdot \psi_n &= 0. \label{eq:Qpsi_free}
\end{align}
It extends to all of $\A_{\text{free}}$ by the graded Leibniz rule $Q(a\cdot b) = (Qa)\cdot b + (-1)^{|a|}a\cdot (Qb)$.
\end{proposition}

\begin{proof}
The BV bracket pairs the field~$\phi$ with the antifield~$\psi$ via $\{\phi(x),\psi(y)\}_{\mathrm{BV}} = \delta^{(3)}(x-y)$. Taking the functional derivative of $S_{\text{free}}$ with respect to the antifield~$\psi$ yields
\[
Q\phi = \frac{\delta S_{\text{free}}}{\delta \psi} = (d_t + \bar{\partial}_z)\phi.
\]
Passing to holomorphic modes $\phi_n = \frac{1}{n!}\partial_z^n\phi$ gives $Q\phi_n = \psi_n$. Since $\psi$ is the BV antifield (of top cohomological degree in the two-term complex), $Q\psi_n = 0$. One verifies directly that
\[
Q^2\phi_n = Q(\psi_n) = 0,
\]
confirming that $(\A_{\text{free}}, Q)$ is a cochain complex. The Leibniz rule follows from the derivation property of the BV bracket.
\end{proof}

\begin{remark}[Free $\neq$ Trivial Differential]
Even though the theory is \emph{free} (no interaction vertices), the differential $Q$ is not zero. It encodes the classical equations of motion via the BV antifield pairing. ``Free theory'' therefore means $m_{k\ge 3}=0$, not $Q=0$.
\end{remark}

\subsubsection{The $\lambda$-Bracket $m_2$}

The binary operation $m_2$ is governed by the free propagator.

\begin{construction}[Free propagator on $\R \times \C$]
\label{const:free_propagator_detailed}
The Green's function for $(d_t + \bar{\partial}_z)$ is
\[
K(t-t',z-z') = \frac{1}{2\pi} \frac{\Theta(t-t')}{z-z'},
\]
with $\Theta$ the Heaviside step function. It produces the unique distributional solution of $(d_t+\bar{\partial}_z)K = \delta(t-t')\delta^{(2)}(z-z')$.
\end{construction}

\begin{proposition}[Explicit $\lambda$-brackets]
\label{prop:free_lambda_brackets_detailed}
The only nonvanishing $\lambda$-bracket is
\begin{align}
\label{eq:free_lambda_phi_psi}
\{\phi_n \,{}_\lambda\, \psi_m\} &= \binom{n+m}{n} (-1)^n \lambda^{n+m},
\end{align}
with all others obtained from sesquilinearity and the $Q$-derivation property.
\end{proposition}

\begin{proof}
Contracting with the propagator yields
\[
\{\phi(z_1)\,{}_\lambda\,\psi(z_2)\} = \frac{\Theta(t_1-t_2)}{2\pi(z_1-z_2)}.
\]
Expanding $e^{\lambda(z_1-z_2)}$ and differentiating $n$ and $m$ times in $z_1,z_2$ respectively gives \eqref{eq:free_lambda_phi_psi}. Sesquilinearity then determines all other brackets.
\end{proof}

\begin{remark}[Sesquilinearity check]
Equation \eqref{eq:free_lambda_phi_psi} immediately implies
\[
\{\partial \phi_n \,{}_\lambda\, \psi_m\} = -\lambda \{\phi_n \,{}_\lambda\, \psi_m\},\qquad
\{\phi_n \,{}_\lambda\, \partial \psi_m\} = (\partial+\lambda)\{\phi_n \,{}_\lambda\, \psi_m\},
\]
verifying the PVA sesquilinearity relations.
\end{remark}

\subsubsection{$m_{k \geq 3}$ Vanish at Chain Level}

\begin{proposition}[No Higher Operations in Free Theories]
\label{prop:free_mk_zero}
For the free chiral multiplet,
\[
m_k = 0 \quad \text{for all } k \geq 3.
\]
\end{proposition}

\begin{proof}
Higher operations $m_k$ arise from Feynman diagrams with interaction vertices. Since $S_{\text{free}}$ contains no interaction terms (it is quadratic in fields), all Feynman diagrams contributing to $m_{k \geq 3}$ vanish identically. There are no vertices to place at internal points, so no diagrams exist.
\end{proof}

\subsubsection{Cohomology Computation}

\begin{proposition}[Cohomology of the Free Chiral Multiplet]
\label{prop:free_cohomology}
We have $H^0(\A_{\text{free}}, Q) = \C$ and $H^k(\A_{\text{free}}, Q) = 0$ for $k \neq 0$.
\end{proposition}

\begin{proof}
Each pair $(\phi_n,\psi_n)$ forms an acyclic $Q$-doublet
\[
\phi_n \xrightarrow{Q} \psi_n \xrightarrow{Q} 0,
\]
so the only nontrivial cohomology class is represented by the constant~$1$.
\end{proof}

\begin{remark}[Trivial Descended PVA]
Since $H^\bullet(\A_{\text{free}}, Q) = \C$, the induced Poisson vertex algebra on cohomology is one-dimensional (the trivial PVA).
\end{remark}

\subsection{Example II: Landau--Ginzburg Model with Cubic Superpotential}
\label{ex:LG}

Landau--Ginzburg (LG) models with polynomial superpotentials provide the first nontrivial examples where $m_3 \neq 0$.

\subsubsection{Setup: LG Theory with $W(\phi) = \lambda_3 \phi^3/3$}

Consider a chiral multiplet $(\phi, \psi)$ with superpotential
\[
W(\phi) = \frac{\lambda_3}{3} \phi^3.
\]
The BV action becomes
\[
S_{\text{LG}} = S_{\text{free}} + \int_{\R \times \C} W'(\phi) \psi \, dt \, dz \, d\bar{z} = S_{\text{free}} + \int_{\R \times \C} \lambda_3 \phi^2 \psi \, dt \, dz \, d\bar{z}.
\]
The interaction term $\phi^2 \psi$ is the \emph{interaction vertex} that generates higher operations.

\subsubsection{BV-BRST Differential}

The differential now includes interaction contributions from the superpotential:
\begin{align}
Q \cdot \phi &= \psi + W'(\phi) = \psi + \lambda_3 \phi^2, \\
Q \cdot \psi &= 0.
\end{align}
Nilpotency $Q^2 = 0$ is verified directly: $Q^2\phi = Q(\psi + \lambda_3\phi^2) = Q\psi + 2\lambda_3\phi\,Q\phi = 0 + 2\lambda_3\phi(\psi + \lambda_3\phi^2)$, which vanishes in the BV complex since the terms are $Q$-exact by the Leibniz rule applied to the full BV action (equivalently, $Q^2 = 0$ follows from the classical master equation $\{S_{\text{LG}}, S_{\text{LG}}\}_{\mathrm{BV}} = 0$).

\begin{remark}
More precisely, on the algebra of local observables $\A_{\text{LG}} = \C[\phi_n, \psi_n]$, the differential is $Q\phi_n = \psi_n$ (the linear part, identical to the free theory), and the superpotential deformation enters through the higher operations $m_2, m_3$ rather than through $m_1 = Q$. The formulas above describe the \emph{full} BV-BRST differential on the space of fields; on the local observable algebra, the interaction is encoded in the $A_\infty$ operations.
\end{remark}

\subsubsection{The First Nontrivial Higher Operation: $m_3$}

\begin{proposition}[Cubic $m_3$ from Superpotential]
\label{prop:LG_m3}
For the LG model with cubic $W$, the operation $m_3$ arises from the triangle diagram with three external legs and one internal vertex. Explicitly, for operators $O_1(z_1)$, $O_2(z_2)$, $O_3(z_3)$:
\begin{equation}
\label{eq:LG_m3_formula}
m_3(O_1, O_2, O_3) = \frac{\lambda_3}{(2\pi i)} \oint_{|w-z_1| < \varepsilon} \frac{dw}{(w-z_1)(w-z_2)(w-z_3)} \times (\text{contractions}),
\end{equation}
where $w$ is the internal vertex position, and ``contractions'' denote propagators connecting $w$ to the external points $z_i$.
\end{proposition}

\begin{proof}[Computation]
The triangle diagram has:
\begin{itemize}
\item One internal vertex at position $w \in \C$, $t_w \in \R$, weighted by the interaction $\lambda_3 \phi^2(w) F(w)$;
\item Three propagators connecting $w$ to the external points $z_1, z_2, z_3$.
\end{itemize}

Integrating over $w$ and $t_w$ yields a logarithmic form on $\FM_3(\C)$ with a triple pole at the collision divisor where all three points approach a common limit. The residue computation produces (\ref{eq:LG_m3_formula}) up to numerical factors depending on the Feynman rules.
\end{proof}

\subsubsection{$m_{k \geq 4}$}

\begin{proposition}[Truncation at $m_3$ for Cubic LG]
\label{prop:LG_truncation}
For the cubic Landau--Ginzburg model, higher operations satisfy:
\begin{itemize}
\item $m_2$ is the free propagator $\lambda$-bracket (\ref{eq:free_lambda_bracket});
\item $m_3$ is given by the triangle diagram (\ref{eq:LG_m3_formula});
\item $m_k = 0$ for all $k \geq 4$.
\end{itemize}
\end{proposition}

\begin{proof}[Proof by degree/dimension counting]
We use a \emph{grading argument} based on ghost number and form degree on FM compactifications.

\textbf{Setup:} Each operator $O_i$ has a ghost number $|O_i|$, and each operation $m_k$ has degree $1-k$ (Definition \ref{def:ainfty_chiral}). Feynman diagrams contributing to $m_k$ must:
\begin{itemize}
\item Have $k$ external legs (operators $O_1,\ldots,O_k$);
\item Satisfy total ghost number conservation: output degree $= \sum_i |O_i| + (1-k)$.
\end{itemize}

\textbf{Cubic vertex constraint:} The interaction $\lambda_3 \phi^2 F$ is a \emph{cubic} vertex (degree 3 in fields). Any Feynman diagram for $m_k$ with $k \geq 4$ external legs requires:
\begin{itemize}
\item At least $k-2$ internal vertices (for a tree diagram), or more for loop diagrams;
\item Each internal vertex contributes a factor of $\lambda_3$ and increases the total field count.
\end{itemize}

\textbf{Dimensional analysis:} The form degree on $\FM_k(\C)$ is $k-1$ (as required by locality). For $k \geq 4$:
\begin{itemize}
\item A tree-level diagram with $k-2 \geq 2$ cubic vertices would require at least 6 external fields from the vertices alone, but we only have $k = 4$ external legs---contradiction.
\item Loop diagrams with $k \geq 4$ external legs and cubic vertices produce integrands of form degree $> k-1$, which cannot be integrated over $(k-1)$-cycles in $\FM_k(\C)$ to yield a well-defined operation.
\end{itemize}

\textbf{Conclusion:} The only nonvanishing operations are:
\begin{itemize}
\item $m_1 = Q$ (differential);
\item $m_2$ (free propagator);
\item $m_3$ (single cubic vertex, tree level).
\end{itemize}
For $k \geq 4$, degree/dimension constraints force $m_k = 0$.
\end{proof}

\begin{remark}[General Polynomial Superpotentials]
For a general polynomial superpotential $W(\phi) = \sum_{d=2}^D \lambda_d \phi^d/d$, the first nontrivial higher operation is $m_D$ (from the highest-degree term). Operations $m_k$ for $k > D$ typically vanish by similar degree counting, though loop contributions can introduce subtleties.
\end{remark}

\subsubsection{Cohomology and Descent to PVA}

\begin{proposition}[LG Cohomology]
\label{prop:LG_cohomology}
Since $m_1 = Q$ is the \emph{linear} part of the BV-BRST differential (Remark above), the complex $(\A_{\text{LG}}, Q)$ has the same acyclic doublet structure as the free theory:
\[
H^\bullet(\A_{\text{LG}}, Q) = \C,
\]
concentrated in degree~$0$. The Jacobian ring $\C[\phi]/(W'(\phi)) = \C[\phi]/(\phi^2)$ emerges from the descended PVA structure via $A_\infty$ homotopy transfer: the operations $m_2$ and $m_3$ (which encode the superpotential) induce a nontrivial PVA on cohomology representatives whose structure is controlled by the critical locus of $W$.
\end{proposition}

\noindent The descended PVA structure on $H^\bullet(\A_{\text{LG}}, Q)$ encodes the quantum geometry of the LG model. Higher operations $m_3$ vanish in cohomology (by Proposition \ref{prop:m3_vanish}) despite being nonzero at the chain level.

\subsection{Example III: Abelian Gauge Theory with Chern--Simons Coupling}
\label{ex:abelian_CS}

Our final example illustrates how gauge theories with Chern--Simons terms produce rich boundary chiral algebra structures.

\subsubsection{3d Abelian Gauge Theory on $\R \times \C$}

Consider $U(1)$ gauge theory on $\R \times \C$ with Chern--Simons coupling:
\[
S_{CS} = \frac{k}{4\pi} \int_{\R \times \C} A \wedge dA,
\]
where $A$ is a $U(1)$ gauge connection and $k \in \Z$ is the \emph{Chern--Simons level}.

After HT twist and BV gauge fixing, the theory has:
\begin{itemize}
\item Gauge field $A = A_z dz + A_{\bar{z}} d\bar{z} + A_t dt$;
\item Ghost fields $c$ (ghost), $\bar{c}$ (antighost);
\item Auxiliary fields completing the BV multiplet.
\end{itemize}

\subsubsection{Boundary Conditions and Chiral Algebras}

The richness of gauge theories emerges when we impose \emph{boundary conditions} at spatial infinity or on actual boundaries. For 3d gauge theory on $\R \times \C$, consider imposing a boundary condition at a fixed value of $t$ (e.g., $t = 0$), creating a boundary $\{t=0\} \cong \C$.

\begin{theorem}[Boundary Chiral Algebra for Abelian CS]
\label{thm:CS_boundary}
Assume {\rm (H1)--(H4)}.
For $U(1)$ Chern--Simons theory at level $k$ with Dirichlet boundary conditions on $\{t=0\} \cong \C$, the boundary chiral algebra is the \emph{affine Kac--Moody algebra} $\widehat{\mathfrak{u}(1)}_k$ at level $k$.
\end{theorem}

\begin{proof}
\textbf{Step 1 (Boundary field).}
Imposing Dirichlet boundary conditions $A_{\bar z}|_{t=0} = 0$ on the boundary $\{t=0\} \cong \C$ leaves the holomorphic component $J(z) := A_z(0,z)$ as the unconstrained boundary degree of freedom. By the equations of motion $(\partial_t + \bar\partial)A_z = 0$ in the bulk, $J(z)$ is a holomorphic field on $\C$.

\textbf{Step 2 (OPE from bulk propagator).}
The bulk propagator for abelian CS at level $k$ is
\[
\langle A_z(t_1,z_1)\, A_{\bar z}(t_2,z_2) \rangle = \frac{k}{2\pi} \cdot \frac{\Theta(t_1 - t_2)}{z_1 - z_2}.
\]
To compute $\langle J(z_1) J(z_2) \rangle$, we use the BV antibracket to pair $A_z$ with $A_{\bar z}$, then restrict to the boundary $t_1 = t_2 = 0$. The singular part arises from the coincident-point limit:
\[
\langle J(z_1) J(z_2) \rangle = \lim_{t_1 \to 0^+, t_2 \to 0^-} \frac{k}{2\pi} \cdot \frac{\Theta(t_1 - t_2)}{z_1 - z_2} = \frac{k}{2\pi} \cdot \frac{1}{z_1 - z_2}.
\]
The $\lambda$-bracket is obtained by taking the Fourier transform (residue):
\begin{equation}
\label{eq:U1_OPE}
\{J_\lambda J\} = k\lambda, \qquad \text{equivalently} \quad J(z) J(w) \sim \frac{k}{(z-w)^2},
\end{equation}
which is the defining OPE of $\widehat{\mathfrak{u}(1)}_k$: a rank-1 affine Kac--Moody algebra at level $k$.

\textbf{Step 3 (Higher operations vanish).}
Since abelian CS is a free (Gaussian) theory after gauge-fixing, all connected Feynman diagrams with $\ge 3$ external legs vanish (there are no interaction vertices). Hence $m_k = 0$ for $k \ge 3$, and the boundary algebra is a strict vertex algebra (not merely $A_\infty$).

\textbf{Step 4 (Non-abelian generalization).}
For non-abelian gauge group $G$ with Lie algebra $\mathfrak{g}$, the boundary current $J^a(z) = A_z^a(0,z)$ generates $\widehat{\mathfrak{g}}_k$ with OPE $J^a(z) J^b(w) \sim k\delta^{ab}/(z-w)^2 + f^{ab}_c J^c(w)/(z-w)$, where $f^{ab}_c$ are structure constants. The cubic vertex $\int f^{ab}_c A^a A^b A^c$ generates non-trivial $m_k$ for $k \ge 3$ (non-abelian corrections); see \cite{CDG20,KZ25} for complete details.
\end{proof}

\subsubsection{Bulk PVA and Bulk--Boundary Correspondence}

\begin{proposition}[Bulk PVA for CS Theory]
Under Hypotheses~{\rm (H1)--(H4)}, the bulk local operators in cohomology $H^\bullet(\A_{\text{bulk}}, Q)$ form a \emph{commutative} Poisson vertex algebra (the commutativity follows from topological invariance in $t$). The $\lambda$-bracket encodes the classical Poisson structure from the Chern--Simons action.
\end{proposition}

The bulk--boundary correspondence (Theorem \ref{thm:bulk_hochschild}) relates bulk observables to chiral Hochschild cochains of the boundary algebra $\widehat{\mathfrak{u}(1)}_k$. This provides a concrete realization of the abstract framework developed in Sections \ref{sec:chiral_hochschild_expanded}--\ref{sec:bar_cobar}.

\begin{remark}[Connection to Quantum Groups]
The boundary affine Kac--Moody algebra $\widehat{\mathfrak{g}}_k$ is intimately related to quantum groups $U_q(\widehat{\mathfrak{g}})$ at $q = e^{2\pi i/(k + h^\vee)}$, where $h^\vee$ is the dual Coxeter number. Line operators in the 3d theory correspond to representations of $U_q(\widehat{\mathfrak{g}})$, providing a field-theoretic origin for quantum group structures.
\end{remark}

\subsection{Summary Table}

\begin{table}[h]
\centering
\begin{tabular}{|l|c|c|c|}
\hline
\textbf{Property} & \textbf{Free Chiral} & \textbf{LG Cubic} & \textbf{Abelian CS} \\
\hline
Fields & $(\phi, \psi)$ & $(\phi, \psi)$ & $(A_z, c, \bar{c})$ \\
\hline
Interaction & None & $\lambda_3 \phi^3/3$ & $k\,A\wedge dA$ \\
\hline
$Q = m_1$ & $Q\phi=\psi$, $Q\psi=0$ & $Q\phi=\psi$, $Q\psi=0$ & Gauge BRST \\
\hline
$m_2$ & Free propagator & Free propagator & Gauge propagator \\
\hline
$m_3$ & $0$ & Triangle diagram & Gauge vertex \\
\hline
$m_{k\ge 4}$ & $0$ & $0$ (degree constraints) & Higher gauge ops \\
\hline
$H^\bullet(\A,Q)$ & $\C$ & $\C$ & Commutative PVA \\
\hline
Boundary algebra & -- & Matrix factorizations & $\widehat{\mathfrak{u}(1)}_k$ \\
\hline
Quantum structure & -- & Jacobian ring & $U_q(\widehat{\mathfrak{u}(1)})$ \\
\hline
\end{tabular}
\caption{Comparison of the three illustrative examples.}
\label{tab:examples_summary}
\end{table}
